\section{Historical Experiments}

\subsection{Daily evaluation}

The common daily scoring window contains 248 forecasts and 37 positives. Table
\ref{tab:daily} reports the historical comparison. M2 leads the initial four models, but rolling
30- and 60-day rates both outperform it after stronger baselines are added. Historical M2 is
therefore a model-development result, not the strongest available reference.

\begin{table}[h]
\centering
\caption{Daily expanding-window results on a common scoring window.}
\label{tab:daily}
\begin{tabular}{lrrr}
\toprule
Model & Brier & Log loss & Brier skill vs.\ M0 \\
\midrule
M0 & 0.129385 & 0.435302 & 0.000\% \\
M1 & 0.127290 & 0.430233 & 1.620\% \\
M2 & 0.125443 & 0.425511 & 3.047\% \\
M3-lite & 0.127800 & 0.434866 & 1.226\% \\
\midrule
Rolling 60d & 0.122391 & 0.411335 & 5.406\% \\
Rolling 30d & \textbf{0.119664} & \textbf{0.394584} & \textbf{7.513\%} \\
M2 without regime & 0.127477 & 0.433266 & 1.475\% \\
\bottomrule
\end{tabular}
\end{table}

\paragraph{Scale and calibration.}
The daily event prevalence is 14.919\%, so it is not extremely rare. Absolute daily Brier values
are interpreted with prevalence, skill, calibration, and PR-AUC. Rolling 30-day predictions have a
calibration intercept of $-0.119$ and slope of $0.935$, whereas M2 has intercept $-0.736$ and slope
$0.470$.

\subsection{Six-hour evaluation}

The six-hour dataset has 1,260 periods and exactly 41 positive periods, matching the announcement
count. Once the fixed 30-day training requirement and two-positive minimum are applied, 989
forecasts with 38 positives remain. M2 again gives the best point estimates
(Table~\ref{tab:sixhour}), but its paired weekly-block Brier difference from M0 has a 95\% interval
of $[-0.002073,0.000235]$. The interval includes zero.

\begin{table}[h]
\centering
\caption{Six-hour expanding-window results. PR-AUC is secondary.}
\label{tab:sixhour}
\begin{tabular}{lrrr}
\toprule
Model & Brier & Log loss & PR-AUC \\
\midrule
M0 & 0.037151 & 0.167388 & 0.077297 \\
M1 & 0.037038 & 0.166331 & 0.068535 \\
M2 & \textbf{0.036443} & \textbf{0.159763} & \textbf{0.085141} \\
M3-lite & 0.037029 & 0.167232 & 0.069052 \\
\bottomrule
\end{tabular}
\end{table}

Splitting incident age into 0--6, 6--24, 24--48, and 48--72 hour bins does not improve the
renewal baseline. Removing incident strength gives the best context ablation, but its Brier skill
over M1 is only 0.080\%. These post-development ablations are diagnostic and do not modify the
frozen prospective models.

\subsection{Bootstrap sensitivity}

Paired block intervals are recomputed with 7-, 14-, and 21-day blocks. M2 minus M0 crosses zero for
every block length, as does M3-lite minus M2. M2 is worse than rolling 60 days in point estimate, and
its difference interval also crosses zero for all three block sizes. With one roughly 315-day
sequence, these intervals are sensitivity summaries rather than stable repeated-sample uncertainty.
